% ============================================================================
% CAPÍTULO 5: RESULTADOS
% ============================================================================

\chapter{Resultados}
\label{cap:resultados}

% ----------------------------------------------------------------------------
% INTRODUCCIÓN
% ----------------------------------------------------------------------------
En este capítulo se presentan los resultados obtenidos tras aplicar la metodología descrita en el \Cref{cap:metodologia}. Los hallazgos se organizan según los objetivos específicos planteados.

% ----------------------------------------------------------------------------
% SECCIÓN 5.1: Análisis Exploratorio
% ----------------------------------------------------------------------------
\section{Análisis Exploratorio de Datos}
\label{sec:exploratorio}

[Presenta estadísticas descriptivas y visualizaciones iniciales]

\subsection{Distribución de Variables}

La \Cref{fig:distribucion} muestra la distribución de las principales variables del dataset.

\begin{figure}[htbp]
    \centering
    \includegraphics[width=0.8\textwidth]{figuras/distribucion.png}
    \caption{Distribución de variables principales. (a) Variable 1, (b) Variable 2, (c) Variable 3}
    \label{fig:distribucion}
\end{figure}

% ----------------------------------------------------------------------------
% SECCIÓN 5.2: Resultados del Modelo
% ----------------------------------------------------------------------------
\section{Desempeño del Modelo Propuesto}
\label{sec:desempeño}

\subsection{Métricas de Evaluación}

La \Cref{tab:resultados} presenta las métricas obtenidas en el conjunto de prueba.

\begin{table}[htbp]
    \centering
    \caption{Resultados del modelo propuesto}
    \label{tab:resultados}
    \begin{tabular}{@{}lcccc@{}}
        \toprule
        \textbf{Clase} & \textbf{Precision} & \textbf{Recall} & \textbf{F1-Score} & \textbf{Soporte} \\
        \midrule
        Clase A        & 0.92               & 0.89            & 0.90              & 450 \\
        Clase B        & 0.88               & 0.91            & 0.89              & 380 \\
        Clase C        & 0.85               & 0.83            & 0.84              & 320 \\
        \midrule
        \textbf{Macro Avg}   & \textbf{0.88}    & \textbf{0.88}   & \textbf{0.88}     & \textbf{1150} \\
        \textbf{Weighted Avg}& \textbf{0.89}    & \textbf{0.88}   & \textbf{0.88}     & \textbf{1150} \\
        \bottomrule
    \end{tabular}
\end{table}

El modelo alcanza un F1-Score promedio de \textbf{0.88}, superando el baseline de...

\subsection{Matriz de Confusión}

La \Cref{fig:confusion} presenta la matriz de confusión del modelo.

\begin{figure}[htbp]
    \centering
    \includegraphics[width=0.6\textwidth]{figuras/confusion_matrix.png}
    \caption{Matriz de confusión del modelo propuesto en conjunto de prueba}
    \label{fig:confusion}
\end{figure}

Se observa que la mayor confusión ocurre entre las clases B y C...

% ----------------------------------------------------------------------------
% SECCIÓN 5.3: Comparación con Baselines
% ----------------------------------------------------------------------------
\section{Comparación con Métodos Baseline}
\label{sec:comparacion}

Se comparó el modelo propuesto con tres métodos baseline. La \Cref{tab:comparacion_metodos} resume los resultados.

\begin{table}[htbp]
    \centering
    \caption{Comparación de métodos}
    \label{tab:comparacion_metodos}
    \begin{tabular}{@{}lcccc@{}}
        \toprule
        \textbf{Método} & \textbf{Accuracy} & \textbf{F1-Score} & \textbf{Tiempo (s)} & \textbf{Parámetros} \\
        \midrule
        Logistic Regression & 0.78  & 0.76  & 0.5   & 125 \\
        Random Forest       & 0.82  & 0.81  & 2.3   & --- \\
        SVM (RBF)          & 0.85  & 0.84  & 5.7   & 1,250 \\
        \textbf{Modelo Propuesto} & \textbf{0.89} & \textbf{0.88} & \textbf{3.2} & \textbf{15,432} \\
        \bottomrule
    \end{tabular}
\end{table}

La \Cref{fig:comparacion_barras} visualiza esta comparación.

\begin{figure}[htbp]
    \centering
    \includegraphics[width=0.7\textwidth]{figuras/comparacion_metodos.png}
    \caption{Comparación de F1-Score entre métodos. El modelo propuesto supera a todos los baselines}
    \label{fig:comparacion_barras}
\end{figure}

% ----------------------------------------------------------------------------
% SECCIÓN 5.4: Análisis de Sensibilidad
% ----------------------------------------------------------------------------
\section{Análisis de Sensibilidad}
\label{sec:sensibilidad}

Se evaluó la sensibilidad del modelo a la variación de hiperparámetros.

\subsection{Tasa de Aprendizaje}

La \Cref{fig:learning_rate} muestra el impacto de la tasa de aprendizaje en el desempeño.

\begin{figure}[htbp]
    \centering
    \includegraphics[width=0.7\textwidth]{figuras/learning_rate_analysis.png}
    \caption{F1-Score vs. tasa de aprendizaje. El valor óptimo se encuentra en $\alpha = 0.001$}
    \label{fig:learning_rate}
\end{figure}

% ----------------------------------------------------------------------------
% SECCIÓN 5.5: Resultados por Objetivo Específico
% ----------------------------------------------------------------------------
\section{Resultados por Objetivo Específico}
\label{sec:por_objetivo}

\subsection{Objetivo Específico 1}

[Presenta los resultados relacionados con OE1]

Los hallazgos indican que...

\subsection{Objetivo Específico 2}

[Presenta los resultados relacionados con OE2]

Como se observa en la \Cref{fig:oe2_resultado}...

% ============================================================================
% CONSEJOS:
% ============================================================================
% - Presenta resultados OBJETIVOS (números, gráficas, tablas)
% - NO interpretes aquí (guarda interpretaciones para Discusión)
% - Usa visualizaciones claras y profesionales
% - Incluye intervalos de confianza/desviaciones estándar cuando sea relevante
% - Conecta cada resultado con un objetivo específico
% - Compara siempre con baselines o estado del arte
% ============================================================================


